% !PRINTED 570
\weberpage{570}

%<<CONT
\weq{8}{P_\mu(x) = \Pi\left[x - \wp\left(\frac{\nu}{\mu}\right)\right],}

und $P_\mu(x)$ ist vom Grade $m - 1$.

In dem Produkt $P_\mu(x)$ kommt jeder Faktor $x - g$ zweimal vor, außer wenn $g$ einem der $e$ gleich ist, wenn also die Kongruenz (7) für ein von Null verschiedenes $\nu$ erfüllt werden kann.

1) Wenn $\mu$ relativ prim zu 2 ist, also wenn $m$ \gesperrt{ungerade} ist, so hat die Kongruenz (7) nur die Wurzel $\nu = 0$. Dann ist $P_\mu(x)$ ein Quadrat, und wenn $S(x)$ eine Funktion vom Grade $\frac{1}{2}(m - 1)$ ist,

\weq{9}{P_\mu(x) = S^2, \quad m \equiv 1 \ (\mathrm{mod}\ 2).}

Wenn $\mu$ mit 2 einen Teiler gemein hat, sind drei Fälle zu unterscheiden.

2) Wenn $\mu$ durch 2 teilbar ist, so muß $\nu$, wenn es der Kongruenz (7) genügen soll, ein Vielfaches von $\frac{1}{2}\mu$ sein, und da es drei von Null verschiedene nach dem Modul 2 inkongruente Zahlen in $\Omega$ gibt, so hat (7) drei Wurzeln. Es kommen also unter den $g$ die drei Werte $e_1, e_2, e_3$ vor, und wir haben:

\weq{10}{P_\mu(x) = \wp'(u)^2 S^2, \quad m \equiv 0 \ (\mathrm{mod}\ 2),}

worin $S$ eine ganze Funktion vom Grade $\frac{1}{2}m - 2$ ist.

3) Wenn 2 im Körper $\Omega$ in zwei (gleiche oder verschiedene) Primideale zerfällt, also wenn $\Delta \equiv 0 \ (\mathrm{mod}\ 4)$ oder $\equiv 1 \ (\mathrm{mod}\ 8)$ ist, so kann $\mu$ durch einen dieser Primfaktoren teilbar sein, ohne durch 2 teilbar zu sein. Dann hat die Kongruenz (7) nur eine von Null verschiedene Wurzel und es kommt unter den Zahlen $g$ nur eine der drei Zahlen (6), etwa $e$, vor. In diesem Falle ist

\weq{11}{P_u(x) = [\wp(u) - e]\, S^2,} % [sic] Das Original druckt P_u; die benachbarten Formeln und der Text erfordern P_\mu.

worin $S$ eine Funktion vom Grade $\frac{1}{2}m - 1$ ist.

In bezug auf die beiden besonderen Fälle $g_3 = 0$ und $g_2 = 0$ ist noch folgendes zu bemerken.

% [sic] Das Original setzt hier j(omega)=0. Nach (3), § 150 ist bei g_3=0 vielmehr j=27*64; die Werte in 4) und 5) sind vertauscht und bleiben originalgetreu erhalten.
4) Wenn $g_3 = 0$ ist, so ist $\Delta = -4$, $j(\omega) = 0$, und es ergibt sich aus der Differentialgleichung §\,46, (18):

\[
\wp'(u)^2 = 4\wp(u)^3 - g_2 \wp(u),
\]

durch die mit Rücksicht auf das Anfangsglied der Entwickelung die $\wp$-Funktion eindeutig bestimmt ist.
